\chapter{Exceptions Personnalisées}
\sloppy

\section{Exercice 1 -- Compte Bancaire}

\subsection{Objectif}

Créer une classe \texttt{CompteBancaire} avec \texttt{retirer()} et \texttt{verser()},
et lever une \texttt{SoldeInsuffisantException} si le montant dépasse le solde.

\subsection{Code source}

\begin{lstlisting}[language=Java, caption={exo1\_2.java -- SoldeInsuffisantException}]
class SoldeInsuffisantException extends Exception {}

public class exo1_2 {

    String code;
    double solde;

    public exo1_2(String code, double solde) {
        this.code = code;
        this.solde = solde;
    }

    void verser(exo1_2 e, double montant) {
        e.solde += montant;
        System.out.println("Nouveau solde : " + e.solde);
    }

    void retirer(exo1_2 e, double montant) throws SoldeInsuffisantException {
        if (montant > e.solde) {
            throw new SoldeInsuffisantException();
        }
        e.solde -= montant;
        System.out.println("Nouveau solde : " + e.solde);
    }
}
\end{lstlisting}

\subsection{Explication}

\texttt{SoldeInsuffisantException} hérite de \texttt{Exception} (checked).
Le mot-clé \texttt{throws} dans la signature de \texttt{retirer()} oblige
l'appelant à gérer cette situation.

\section{Exercice 2 -- Montant Invalide}

\subsection{Objectif}

Créer une \texttt{MontantInvalideException} levée lorsque le montant est $\leq 0$.

\subsection{Code source}

\begin{lstlisting}[language=Java, caption={exo2\_2.java -- MontantInvalideException}]
class MontantInvalideException extends RuntimeException {}

public class exo2_2 {

    String code;
    double solde;

    public exo2_2(String code, double solde) {
        this.code = code;
        this.solde = solde;
    }

    void verser(exo2_2 e, double montant) {
        if (montant <= 0) {
            throw new MontantInvalideException();
        }
        e.solde += montant;
        System.out.println("Nouveau solde : " + e.solde);
    }
}
\end{lstlisting}

\subsection{Explication}

\texttt{MontantInvalideException} hérite de \texttt{RuntimeException}
(unchecked), car un montant invalide est une erreur de logique qui ne devrait
jamais se produire en production correcte.

\section{Exercice 3 -- Double Validation}

\subsection{Objectif}

Modifier \texttt{retirer()} pour gérer à la fois \texttt{MontantInvalideException}
et \texttt{SoldeInsuffisantException}.

\subsection{Code source}

\begin{lstlisting}[language=Java, caption={exo3\_2.java -- Double validation}]
class SoldeInsuffisantException3 extends Exception {}
class MontantInvalideException3 extends RuntimeException {}

public class exo3_2 {

    String code;
    double solde;

    public exo3_2(String code, double solde) {
        this.code = code;
        this.solde = solde;
    }

    void retirer(exo3_2 e, double montant) throws SoldeInsuffisantException3 {
        if (montant <= 0) {
            throw new MontantInvalideException3();
        }
        if (montant > e.solde) {
            throw new SoldeInsuffisantException3();
        }
        e.solde -= montant;
        System.out.println("Nouveau solde : " + e.solde);
    }
}
\end{lstlisting}

\section{Exercice 4 -- Inscription Utilisateur}

\subsection{Objectif}

Valider l'email et l'âge lors d'une inscription via deux exceptions personnalisées.

\subsection{Code source}

\begin{lstlisting}[language=Java, caption={exo4\_2.java -- EmailInvalideException \& AgeInvalideException}]
class EmailInvalideException4 extends Exception {}
class AgeInvalideException4 extends Exception {}

public class exo4_2 {

    void inscrire(exo4_2 e, String email, int age)
            throws EmailInvalideException4, AgeInvalideException4 {
        if (!email.contains("@")) {
            throw new EmailInvalideException4();
        }
        if (age < 18) {
            throw new AgeInvalideException4();
        }
        System.out.println("Inscription reussie");
    }
}
\end{lstlisting}

\section{Exercice 5 -- Authentification}

\subsection{Objectif}

Lever une \texttt{AuthentificationException} avec le message
\og Identifiants incorrects \fg{} en cas de mauvais login.

\subsection{Code source}

\begin{lstlisting}[language=Java, caption={exo5\_2.java -- AuthentificationException}]
class AuthentificationException5 extends Exception {}

public class exo5_2 {

    void login(exo5_2 e, String username, String password)
            throws AuthentificationException5 {
        if (!username.equals("admin") || !password.equals("1234")) {
            throw new AuthentificationException5();
        }
        System.out.println("Connexion reussie");
    }

    public static void main(String[] args) {
        exo5_2 e = new exo5_2();
        try {
            e.login(e, "user", "wrong");
        } catch (AuthentificationException5 ex) {
            System.out.println("Identifiants incorrects");
        }
    }
}
\end{lstlisting}

\section{Exercice 6 -- Gestion de Stock}

\subsection{Objectif}

Lever une \texttt{StockInsuffisantException} si la quantité demandée dépasse le stock disponible.

\subsection{Code source}

\begin{lstlisting}[language=Java, caption={exo6\_2.java -- StockInsuffisantException}]
class StockInsuffisantException6 extends Exception {}

public class exo6_2 {
    int stock;

    public exo6_2(int stock) {
        this.stock = stock;
    }

    void retirerDuStock(exo6_2 p, int quantite)
            throws StockInsuffisantException6 {
        if (quantite > p.stock) {
            throw new StockInsuffisantException6();
        }
        p.stock -= quantite;
        System.out.println("Stock restant : " + p.stock);
    }
}
\end{lstlisting}

\section{Exercice 7 -- Téléchargement}

\subsection{Objectif}

Lever une \texttt{QuotaDepasseException} si la taille d'un fichier dépasse la limite autorisée.

\subsection{Code source}

\begin{lstlisting}[language=Java, caption={exo7\_2.java -- QuotaDepasseException}]
class QuotaDepasseException7 extends Exception {}

public class exo7_2 {
    double limite = 500.0;

    void telechargerFichier(exo7_2 e, double taille)
            throws QuotaDepasseException7 {
        if (taille > e.limite) {
            throw new QuotaDepasseException7();
        }
        System.out.println("Telechargement fini");
    }
}
\end{lstlisting}

\section{Exercice 8 -- Validation de Formulaire}

\subsection{Objectif}

Lever une \texttt{ChampObligatoireException} si un champ du formulaire est vide ou nul.

\subsection{Code source}

\begin{lstlisting}[language=Java, caption={exo8\_2.java -- ChampObligatoireException}]
class ChampObligatoireException8 extends Exception {}

public class exo8_2 {

    void validerFormulaire(exo8_2 e, String nom, String email)
            throws ChampObligatoireException8 {
        if (nom == null || nom.isEmpty() || email == null || email.isEmpty()) {
            throw new ChampObligatoireException8();
        }
        System.out.println("Formulaire valide");
    }
}
\end{lstlisting}

\section{Exercice 9 -- Paiement}

\subsection{Objectif}

Gérer deux exceptions lors d'un paiement~: dépassement de plafond et carte expirée.

\subsection{Code source}

\begin{lstlisting}[language=Java, caption={exo9\_2.java -- PaiementRefuseException \& CarteExpireeException}]
class PaiementRefuseException9 extends Exception {}
class CarteExpireeException9 extends Exception {}

public class exo9_2 {
    double plafond = 2000.0;

    void payer(exo9_2 e, double montant, boolean expiree)
            throws PaiementRefuseException9, CarteExpireeException9 {
        if (expiree) {
            throw new CarteExpireeException9();
        }
        if (montant > e.plafond) {
            throw new PaiementRefuseException9();
        }
        System.out.println("Paiement effectue");
    }
}
\end{lstlisting}

\section{Exercice 10 -- Conception~: Quand Utiliser Chaque Type ?}

\subsection{Objectif}

Répondre aux questions de conception sur le choix du type d'exception.

\subsection{Synthèse}

\begin{table}[H]
\centering
\begin{tabularx}{0.95\textwidth}{L{4cm} X}
\toprule
\textbf{\color{primarycolor}Question} &
\textbf{\color{primarycolor}Réponse} \\
\midrule
Quand utiliser \textit{checked} ? &
Pour les erreurs \textbf{récupérables} et externes~: fichier introuvable,
connexion réseau, base de données. Le compilateur force la gestion. \\
\addlinespace[0.2cm]
Quand utiliser \textit{unchecked} ? &
Pour les \textbf{bugs logiques} du développeur~: index hors limites,
argument nul, montant invalide. Pas de \texttt{throws} requis. \\
\addlinespace[0.2cm]
Pourquoi créer une exception personnalisée ? &
Pour donner un \textbf{sens métier} précis à l'erreur, améliorer la
lisibilité du code et faciliter la maintenance. \\
\bottomrule
\end{tabularx}
\caption{Synthèse sur la conception des exceptions}
\end{table}
